\documentclass[11pt,a4paper,twoside]{article}

% ============================================================================
% PAQUETES ESENCIALES
% ============================================================================
\usepackage[spanish]{babel}
\usepackage[margin=2.5cm]{geometry}
\usepackage{graphicx}
\usepackage{booktabs}
\usepackage{array}
\usepackage{multirow}
\usepackage{hyperref}
\usepackage{xcolor}
\usepackage{fancyhdr}
\usepackage{lastpage}
\usepackage{float}
\usepackage{amsmath}
\usepackage{amssymb}
\usepackage{tikz}
\usepackage{pgfplots}
\usepackage{listings}
\usepackage{xfrac}
\usepackage{caption}
\usepackage{subcaption}

% ============================================================================
% CONFIGURACIÓN DE COLORES Y ESTILOS
% ============================================================================
\definecolor{darkblue}{rgb}{0.0, 0.2, 0.6}
\definecolor{lightblue}{rgb}{0.0, 0.4, 0.8}
\definecolor{accent}{rgb}{1.0, 0.3, 0.0}
\definecolor{darkred}{rgb}{0.8, 0.0, 0.0}

\hypersetup{
    colorlinks=true,
    linkcolor=darkblue,
    urlcolor=lightblue,
    citecolor=darkblue
}

% ============================================================================
% ENCABEZADO Y PIE DE PÁGINA
% ============================================================================
\pagestyle{fancy}
\fancyhf{}
\fancyhead[LE,RO]{\textcolor{darkblue}{\textbf{People Analytics: FinanCorp Chile}}}
\fancyhead[LO,RE]{\small\textit{Diagnóstico de Rotación y Retención de Talento}}
\fancyfoot[LE,RO]{\thepage\ de \pageref{LastPage}}
\fancyfoot[LO,RE]{\small\textit{Mayo 2026}}
\renewcommand{\headrulewidth}{1.5pt}
\renewcommand{\footrulewidth}{0.5pt}

% ============================================================================
% CONFIGURACIÓN DE TÍTULOS Y SECCIONES
% ============================================================================
\usepackage{titlesec}

\titleformat{\section}
  {\normalfont\Large\bfseries\color{darkblue}}
  {\thesection}{1em}{}[\titlerule]

\titleformat{\subsection}
  {\normalfont\large\bfseries\color{lightblue}}
  {\thesubsection}{0.8em}{}

\titleformat{\subsubsection}
  {\normalfont\normalsize\bfseries\color{accent}}
  {\thesubsubsection}{0.6em}{}

% ============================================================================
% CONFIGURACIÓN DE FIGURAS Y TABLAS
% ============================================================================
\captionsetup{font=small, labelfont=bf, textfont=it}
\setlength{\parindent}{0em}
\setlength{\parskip}{0.5em}

% ============================================================================
% INICIO DEL DOCUMENTO
% ============================================================================
\begin{document}

% ============================================================================
% PORTADA
% ============================================================================
\thispagestyle{empty}
\begin{center}
    \vspace*{2cm}
    
    {\fontsize{28}{34}\selectfont\color{darkblue}\textbf{PEOPLE ANALYTICS}}
    
    \vspace{0.5cm}
    
    {\fontsize{18}{22}\selectfont\color{lightblue}\textbf{Diagnóstico de Rotación, Retención y Fuga de Talento}}
    
    \vspace{0.5cm}
    
    {\fontsize{14}{17}\selectfont\color{accent}\textit{FinanCorp Chile — Caso de Estudio}}
    
    \vspace{3cm}
    
    \begin{tikzpicture}
        \draw[darkblue, line width=0.5mm] (0,0) -- (15,0);
        \draw[lightblue, line width=0.3mm] (0,-0.3) -- (15,-0.3);
    \end{tikzpicture}
    
    \vspace{3cm}
    
    \begin{tabular}{ll}
        \textbf{Empresa:} & FinanCorp Chile (Caso Ficticio) \\
        \textbf{Dotación:} & 3,200 colaboradores \\
        \textbf{Rotación Actual:} & 22\% (\(\sim\)704 salidas anuales) \\
        \textbf{Costo de Rotación:} & USD \$2.112.000 anuales \\
        \textbf{Fecha de Análisis:} & 13 de mayo de 2026
    \end{tabular}
    
    \vfill
    
    {\large\textbf{Análisis Predictivo basado en Machine Learning}}
    
    {\small\textit{Modelo Random Forest — ROC-AUC: 82\%}}
    
    \vspace{2cm}
\end{center}

\newpage

% ============================================================================
% RESUMEN EJECUTIVO
% ============================================================================
\section*{Resumen Ejecutivo}

Este análisis implementa un modelo analítico avanzado para identificar, explicar y anticipar la fuga de talento en organizaciones del sector financiero. Mediante técnicas de análisis exploratorio de datos, segmentación de riesgo y modelado predictivo con Machine Learning, se proporciona una base data-driven para decisiones estratégicas de retención de personal.

\subsection*{Hallazgos Clave}

\begin{itemize}
    \item \textbf{Rotación Crítica:} La rotación en áreas críticas (28.1\%) es \textbf{3.5 veces superior} a otras áreas (7.9\%)
    \item \textbf{Impacto Financiero:} Pérdida anual de USD \$2.112.000 en costos de reemplazo
    \item \textbf{Oportunidad:} Reducir rotación 15\% implicaría ahorrar USD \$525.000 anuales
    \item \textbf{Predicción Efectiva:} Modelo con 82\% de precisión predictiva (ROC-AUC)
    \item \textbf{Intervención Selectiva:} 234 personas en riesgo muy alto requieren atención inmediata
\end{itemize}

\subsection*{Recomendación Principal}

Implementar un programa de intervención en tres niveles con enfoque en las 234 personas clasificadas en riesgo muy alto, proyectando un impacto de USD \$1.356.000 en ahorro directo durante el primer año.

\newpage

% ============================================================================
% TABLA DE CONTENIDOS
% ============================================================================
\tableofcontents
\newpage

% ============================================================================
% SECCIÓN 1: INTRODUCCIÓN
% ============================================================================
\section{Introducción}

\subsection{Contexto y Problemática}

La rotación de personal representa uno de los mayores desafíos estratégicos para las organizaciones modernas. En el sector financiero, donde la continuidad operacional y la expertise son críticos, la fuga de talento genera impactos significativos en:

\begin{itemize}
    \item \textbf{Costos Directos:} Reclutamiento, selección y capacitación de nuevos colaboradores
    \item \textbf{Costos Indirectos:} Pérdida de productividad, disrupciones operacionales, impacto en clima
    \item \textbf{Riesgos Estratégicos:} Vulnerabilidad en áreas críticas, pérdida de conocimiento organizacional
    \item \textbf{Reputacionales:} Deterioro de marca empleadora, dificultad en atraer talento
\end{itemize}

FinanCorp Chile enfrenta una rotación del 22\% anual, significativamente por encima de la industria (15\%), generando una pérdida estimada de USD \$2.112.000 anuales. Esta situación requiere una intervención estratégica basada en datos.

\subsection{Objetivos del Análisis}

\begin{enumerate}
    \item Construir un modelo predictivo para identificar colaboradores en riesgo de rotación
    \item Explicar los factores determinantes en la decisión de abandono
    \item Priorizar intervenciones por impacto y urgencia
    \item Proporcionar recomendaciones estratégicas con ROI estimado
\end{enumerate}

\subsection{Metodología}

Este análisis combina técnicas de:

\begin{itemize}
    \item \textbf{Análisis Exploratorio:} Estadísticas descriptivas y correlaciones
    \item \textbf{Segmentación:} Clasificación de riesgo mediante scoring
    \item \textbf{Machine Learning:} Comparación de Regresión Logística vs. Random Forest
    \item \textbf{Análisis de Impacto:} Matriz de riesgo-urgencia y priorización
\end{itemize}

\newpage

% ============================================================================
% SECCIÓN 2: DESCRIPCIÓN DEL DATASET
% ============================================================================
\section{Datos y Metodología}

\subsection{Características del Dataset}

Se analizó una muestra de 3,200 colaboradores con las siguientes variables:

\begin{table}[H]
    \centering
    \caption{Variables del Análisis}
    \begin{tabular}{|l|l|p{6cm}|}
        \toprule
        \textbf{Variable} & \textbf{Tipo} & \textbf{Descripción} \\
        \midrule
        Antigüedad (años) & Continua & Tiempo en la organización \\
        Desempeño (1-5) & Discreta & Evaluación de rendimiento \\
        Clima (1-5) & Discreta & Percepción del ambiente laboral \\
        Compromiso (1-5) & Discreta & Nivel de engagement y lealtad \\
        Ausentismo (días/año) & Continua & Inasistencias anuales \\
        Movilidad Interna & Binaria & Traslados en últimos 2 años \\
        Capacitación (cursos/año) & Continua & Inversión en desarrollo \\
        Área Crítica & Binaria & Pertenecer a área crítica (sí/no) \\
        Salario (USD) & Continua & Compensación monetaria \\
        \textbf{Rotación} & \textbf{Binaria} & \textbf{Variable objetivo} \\
        \bottomrule
    \end{tabular}
\end{table}

\subsection{Distribución de la Muestra}

La distribución de colaboradores por área refleja una estructura típica de empresa financiera:

\begin{table}[H]
    \centering
    \caption{Distribución de Colaboradores por Área}
    \begin{tabular}{|l|r|r|r|}
        \toprule
        \textbf{Área} & \textbf{Cantidad} & \textbf{\%} & \textbf{Crítica} \\
        \midrule
        Tecnología & 800 & 25.0\% & \checkmark \\
        Operaciones Financieras & 704 & 22.0\% & \checkmark \\
        Atención de Clientes & 640 & 20.0\% & \checkmark \\
        Riesgo y Cumplimiento & 576 & 18.0\% & \checkmark \\
        Finanzas & 256 & 8.0\% & \\
        Personas & 128 & 4.0\% & \\
        Legal & 64 & 2.0\% & \\
        Administración & 32 & 1.0\% & \\
        \midrule
        \textbf{TOTAL} & \textbf{3,200} & \textbf{100\%} & \\
        \bottomrule
    \end{tabular}
\end{table}

\newpage

% ============================================================================
% SECCIÓN 3: INDICADORES CLAVE
% ============================================================================
\section{Indicadores Clave de Rotación}

\subsection{Rotación Global y por Áreas Críticas}

El análisis revela una disparidad crítica en las tasas de rotación:

\begin{table}[H]
    \centering
    \caption{Indicadores de Rotación por Criticidad}
    \begin{tabular}{|l|r|r|r|}
        \toprule
        \textbf{Métrica} & \textbf{Valor} & \textbf{Salidas/Año} & \textbf{Costo (USD)} \\
        \midrule
        Rotación Total & 22.0\% & 704 & \$2.112.000 \\
        Áreas Críticas & 28.1\% & 437 & \$1.311.000 \\
        Otras Áreas & 7.9\% & 267 & \$801.000 \\
        \midrule
        \textbf{Diferencial} & \textbf{20.2 pp} & \textbf{170} & \textbf{\$510.000} \\
        \bottomrule
    \end{tabular}
\end{table}

\subsection{Rotación por Antigüedad}

La antigüedad es un factor protector significativo. Colaboradores con menos de 1 año presentan rotación del 31.2\%, mientras que aquellos con más de 10 años solo 3.8\%.

\subsection{Rotación por Nivel de Desempeño}

Paradójicamente, colaboradores con bajo desempeño presentan menor rotación (19.2\%), sugiriendo que pueden estar presionados a permanencer. El mayor riesgo está en desempeño medio-bueno (25.3\%).

\subsection{Rotación por Capacitación}

Existe correlación inversa significativa: colaboradores sin capacitación tienen 26.1\% de rotación, mientras que quienes reciben 5+ cursos anuales solo 15.3\%.

\newpage

% ============================================================================
% SECCIÓN 4: SEGMENTACIÓN DE RIESGO
% ============================================================================
\section{Segmentación y Análisis de Riesgo}

\subsection{Metodología de Score de Riesgo}

Se desarrolló un modelo de puntuación (0-100) basado en factores predictores:

\[
\text{Score Riesgo} = \sum_{i=1}^{n} w_i \times f_i(x_i)
\]

Donde los pesos asignados son:

\begin{table}[H]
    \centering
    \caption{Pesos de Factores en Score de Riesgo}
    \begin{tabular}{|l|r|l|}
        \toprule
        \textbf{Factor} & \textbf{Peso} & \textbf{Criterio} \\
        \midrule
        Compromiso Bajo & 35\% & Compromiso $<$ 2.5/5 \\
        Ausentismo Alto & 25\% & Ausentismo $>$ 12 días/año \\
        Bajo Desempeño & 20\% & Desempeño $<$ 2.5/5 \\
        Falta Capacitación & 10\% & Cursos/año $<$ 1 \\
        Estancamiento & 5\% & Sin movilidad + 3+ años \\
        Área Crítica & 5\% & Pertenecer a área crítica \\
        \bottomrule
    \end{tabular}
\end{table}

\subsection{Distribución de Riesgo}

\begin{table}[H]
    \centering
    \caption{Distribución por Nivel de Riesgo}
    \begin{tabular}{|l|r|r|r|}
        \toprule
        \textbf{Nivel de Riesgo} & \textbf{Cantidad} & \textbf{\%} & \textbf{Rotación Real} \\
        \midrule
        Bajo (0-25) & 1,186 & 37.1\% & 6.2\% \\
        Medio (25-50) & 1,218 & 38.1\% & 15.8\% \\
        Alto (50-75) & 546 & 17.1\% & 38.5\% \\
        Muy Alto (75-100) & 250 & 7.8\% & 64.8\% \\
        \midrule
        \textbf{TOTAL} & \textbf{3,200} & \textbf{100\%} & \textbf{22.0\%} \\
        \bottomrule
    \end{tabular}
\end{table}

\textbf{Hallazgo Crítico:} Los 250 colaboradores en riesgo muy alto (7.8\%) representan el 72.5\% de todas las rotaciones, con una tasa de 64.8\% versus 22.0\% promedio.

\newpage

% ============================================================================
% SECCIÓN 5: MODELOS PREDICTIVOS
% ============================================================================
\section{Modelo Predictivo de Machine Learning}

\subsection{Comparación de Modelos}

Se entrenaron y compararon dos algoritmos:

\begin{table}[H]
    \centering
    \caption{Comparación de Modelos Predictivos}
    \begin{tabular}{|l|r|r|r|r|}
        \toprule
        \textbf{Modelo} & \textbf{Precisión} & \textbf{ROC-AUC} & \textbf{Sensibilidad} & \textbf{Especificidad} \\
        \midrule
        Regresión Logística & 73.0\% & 0.780 & 68.5\% & 74.2\% \\
        Random Forest & 76.0\% & 0.820 & 71.3\% & 76.8\% \\
        \bottomrule
    \end{tabular}
\end{table}

\textbf{Modelo Seleccionado:} Random Forest (82\% ROC-AUC)

\subsection{Importancia de Features en Random Forest}

El ranking de importancia de variables es:

\begin{enumerate}
    \item \textbf{Compromiso:} 28.4\% — Factor dominante en decisión de rotación
    \item \textbf{Ausentismo:} 22.1\% — Señal de problemas subyacentes
    \item \textbf{Capacitación:} 18.7\% — Inversión en desarrollo
    \item \textbf{Desempeño:} 16.3\% — Evaluación de rendimiento
    \item \textbf{Antigüedad:} 8.2\% — Factor protector
    \item \textbf{Movilidad Interna:} 4.1\% — Oportunidades de carrera
    \item \textbf{Clima:} 1.6\% — Ambiente laboral
    \item \textbf{Área Crítica:} 0.6\% — Criticidad del rol
\end{enumerate}

\subsection{Matriz de Confusión}

\[
\begin{pmatrix}
\text{VN} = 1,931 & \text{FP} = 262 \\
\text{FN} = 117 & \text{VP} = 150
\end{pmatrix}
\]

Donde: VN (Verdaderos Negativos), FP (Falsos Positivos), FN (Falsos Negativos), VP (Verdaderos Positivos)

\newpage

% ============================================================================
% SECCIÓN 6: ANÁLISIS DE ÁREAS CRÍTICAS
% ============================================================================
\section{Análisis Detallado de Áreas Críticas}

\subsection{Tecnología}

\begin{itemize}
    \item \textbf{Total:} 800 colaboradores
    \item \textbf{Rotación:} 30.5\% (244 salidas)
    \item \textbf{Riesgo Alto/Muy Alto:} 18.4\% (147 personas)
    \item \textbf{Problema Principal:} Bajo compromiso (3.1/5) y ausentismo elevado
    \item \textbf{Costo Anual:} \$732.000
\end{itemize}

\subsection{Operaciones Financieras}

\begin{itemize}
    \item \textbf{Total:} 704 colaboradores
    \item \textbf{Rotación:} 28.4\% (200 salidas)
    \item \textbf{Riesgo Alto/Muy Alto:} 16.2\% (114 personas)
    \item \textbf{Problema Principal:} Falta de capacitación y desarrollo
    \item \textbf{Costo Anual:} \$600.000
\end{itemize}

\subsection{Riesgo y Cumplimiento}

\begin{itemize}
    \item \textbf{Total:} 576 colaboradores
    \item \textbf{Rotación:} 25.9\% (149 salidas)
    \item \textbf{Riesgo Alto/Muy Alto:} 14.1\% (81 personas)
    \item \textbf{Problema Principal:} Clima laboral deteriorado (2.8/5)
    \item \textbf{Costo Anual:} \$447.000
\end{itemize}

\subsection{Atención de Clientes}

\begin{itemize}
    \item \textbf{Total:} 640 colaboradores
    \item \textbf{Rotación:} 26.1\% (167 salidas)
    \item \textbf{Riesgo Alto/Muy Alto:} 15.8\% (101 personas)
    \item \textbf{Problema Principal:} Baja movilidad interna y estancamiento
    \item \textbf{Costo Anual:} \$501.000
\end{itemize}

\newpage

% ============================================================================
% SECCIÓN 7: RECOMENDACIONES ESTRATÉGICAS
% ============================================================================
\section{Recomendaciones Estratégicas}

\subsection{Nivel 1: Intervención Urgente (Próximas 2 Semanas)}

\subsubsection{Acción 1: Reuniones 1-a-1 con Riesgos Muy Altos}

\begin{itemize}
    \item \textbf{Población:} 250 colaboradores (7.8\%)
    \item \textbf{Objetivo:} Entender necesidades, identificar pain points
    \item \textbf{Responsable:} Jefaturas directas + Gerencia Personas
    \item \textbf{Duración:} 60 minutos por persona
    \item \textbf{Inversión:} 250 horas (USD \$12.500 estimado)
\end{itemize}

\subsubsection{Acción 2: Diagnóstico Rápido en Áreas Críticas}

\begin{itemize}
    \item Entrevistas de salida profundas (últimos 3 meses)
    \item Encuesta de clima focalizada (48 horas)
    \item Análisis de brecha salarial vs. mercado
\end{itemize}

\subsection{Nivel 2: Intervención Preventiva (Próximo Mes)}

\subsubsection{Acción 1: Programa de Mentoría Estructurada}

\begin{itemize}
    \item Asignar mentores a 546 personas en riesgo medio-alto
    \item Sesiones mensuales de seguimiento (1 hora)
    \item Enfoque en desarrollo de competencias
    \item \textbf{Inversión Estimada:} USD \$85.000/año
\end{itemize}

\subsubsection{Acción 2: Mejora de Clima Laboral}

\begin{itemize}
    \item Auditoría de cargas de trabajo en Tecnología y Riesgo
    \item Fortalecimiento de comunicación de jefaturas (talleres)
    \item Programa de bienestar laboral
\end{itemize}

\subsubsection{Acción 3: Plan de Desarrollo y Capacitación}

\begin{itemize}
    \item Meta: 3+ cursos/año para población en riesgo
    \item Presupuesto estimado: USD \$120.000
    \item Enfoque en liderazgo y competencias técnicas
\end{itemize}

\subsection{Nivel 3: Transformación Organizacional (6 Meses)}

\subsubsection{Acción 1: Dashboard Ejecutivo de Rotación}

\begin{itemize}
    \item KPIs en tiempo real por área
    \item Alertas automáticas por cambios en score de riesgo
    \item Seguimiento mensual de iniciativas
    \item \textbf{Inversión:} USD \$45.000 (desarrollo e implementación)
\end{itemize}

\subsubsection{Acción 2: Integración de People Analytics}

\begin{itemize}
    \item Monthly Reviews de rotación con Jefaturas
    \item Data-driven decisions para promotions/transfers
    \item Presupuesto anual de retención vs. rotación
\end{itemize}

\subsubsection{Acción 3: Revisión de Políticas de Personas}

\begin{itemize}
    \item Análisis de competitividad salarial
    \item Revisión de políticas de carrera
    \item Programa de engagement periódico
\end{itemize}

\newpage

% ============================================================================
% SECCIÓN 8: ANÁLISIS DE IMPACTO
% ============================================================================
\section{Proyección de Impacto Financiero}

\subsection{Escenario Base: Sin Intervención}

\begin{table}[H]
    \centering
    \caption{Escenario Base (Status Quo)}
    \begin{tabular}{|l|r|r|}
        \toprule
        \textbf{Métrica} & \textbf{Año 1} & \textbf{Año 3} \\
        \midrule
        Rotación Global & 22.0\% & 22.0\% \\
        Salidas Anuales & 704 & 704 \\
        Costo Total & \$2.112.000 & \$2.112.000 \\
        \bottomrule
    \end{tabular}
\end{table}

\subsection{Escenario Propuesto: Con Intervención}

\begin{table}[H]
    \centering
    \caption{Proyección con Intervención (Año 1)}
    \begin{tabular}{|l|r|r|r|}
        \toprule
        \textbf{Métrica} & \textbf{Meta} & \textbf{Estimado} & \textbf{Impacto} \\
        \midrule
        Rotación Total & 18.7\% & 598 & -106 salidas \\
        Rotación en Críticas & 22.4\% & 347 & -90 salidas \\
        Ahorro Directo & - & - & \$588.000 \\
        Costo Intervención & - & - & (\$262.500) \\
        \midrule
        \textbf{Beneficio Neto} & - & - & \textbf{\$325.500} \\
        \bottomrule
    \end{tabular}
\end{table}

\subsection{Retorno sobre la Inversión (ROI)}

\[
\text{ROI} = \frac{\text{Beneficio Neto}}{\text{Inversión}} \times 100 = \frac{\$325.500}{\$262.500} \times 100 = 123.9\%
\]

\textbf{Interpretación:} Por cada USD 1 invertido en retención, se esperaría obtener USD 2.24 en beneficios netos (considerando costo de reemplazo evitado).

\subsection{Proyección a 3 Años}

\begin{table}[H]
    \centering
    \caption{Proyección de Ahorros Acumulados (3 Años)}
    \begin{tabular}{|l|r|r|r|}
        \toprule
        \textbf{Concepto} & \textbf{Año 1} & \textbf{Año 2} & \textbf{Año 3} \\
        \midrule
        Ahorro Directo & \$588.000 & \$756.000 & \$840.000 \\
        Inversión & (\$262.500) & (\$150.000) & (\$100.000) \\
        Beneficio Neto & \$325.500 & \$606.000 & \$740.000 \\
        \midrule
        \textbf{Acumulado} & \$325.500 & \$931.500 & \$1.671.500 \\
        \bottomrule
    \end{tabular}
\end{table}

\newpage

% ============================================================================
% SECCIÓN 9: CONCLUSIONES
% ============================================================================
\section{Conclusiones}

\subsection{Hallazgos Principales}

\begin{enumerate}
    \item \textbf{Crisis de Rotación en Áreas Críticas:} La rotación del 28.1\% en áreas críticas versus 7.9\% en otras crea un riesgo de continuidad operacional.
    
    \item \textbf{Factores Predictores Claros:} El compromiso (28.4\%) y ausentismo (22.1\%) son los principales drivers, permitiendo intervención focalizada.
    
    \item \textbf{Modelo Predictivo Efectivo:} El Random Forest con ROC-AUC 0.82 proporciona capacidad fiable para identificación temprana.
    
    \item \textbf{Oportunidad de Intervención Selectiva:} Solo 250 personas (7.8\%) representan el 72.5\% de rotaciones, permitiendo alto impacto con inversión moderada.
    
    \item \textbf{ROI Comprobado:} Intervenciones estratégicas generarían USD \$1.671.500 en beneficios netos a 3 años.
\end{enumerate}

\subsection{Factores Críticos de Éxito}

\begin{itemize}
    \item \textbf{Compromiso Ejecutivo:} Apoyo de C-Suite en implementación
    \item \textbf{Equipos Capacitados:} Jefaturas entrenadas en retención
    \item \textbf{Seguimiento Riguroso:} Monthly reviews de KPIs
    \item \textbf{Inversión Sostenida:} Presupuesto asegurado 3 años
\end{itemize}

\subsection{Riesgos a Considerar}

\begin{itemize}
    \item Resistencia al cambio en jefaturas
    \item Disponibilidad de recursos (HR, presupuesto)
    \item Mercado laboral competitivo externo
    \item Correlación imperfecta entre predictores y rotación real
\end{itemize}

\newpage

% ============================================================================
% SECCIÓN 10: PRÓXIMOS PASOS
% ============================================================================
\section{Plan de Implementación}

\subsection{Semana 1: Validación y Comunicación}

\begin{itemize}
    \item Presentación a Junta Directiva
    \item Obtención de aprobación presupuestaria
    \item Comunicación a liderazgo operativo
\end{itemize}

\subsection{Semana 2: Conformación de Equipos}

\begin{itemize}
    \item Equipo de Gestión de Proyecto (Personas, BI, Finance)
    \item Capacitación de Jefaturas en herramientas
    \item Asignación de responsables por área
\end{itemize}

\subsection{Semana 3: Comenzar Intervención Nivel 1}

\begin{itemize}
    \item Calendarización de reuniones 1-a-1
    \item Iniciación de entrevistas de salida profundas
    \item Encuesta de clima rápida
\end{itemize}

\subsection{Mes 2-3: Intervención Nivel 2 y Dashboard}

\begin{itemize}
    \item Implementación programa de mentoría
    \item Desarrollo de dashboard de rotación
    \item Plan de capacitación personalizado
\end{itemize}

\subsection{Mes 4-6: Transformación Organizacional}

\begin{itemize}
    \item Revisión de políticas de personas
    \item Análisis de competitividad salarial
    \item Integración de People Analytics en decisiones
\end{itemize}

\newpage

% ============================================================================
% BIBLIOGRAFÍA
% ============================================================================
\begin{thebibliography}{9}

\bibitem{pfeffer1999}
Pfeffer, J., \& Veiga, J. F. (1999).
\textit{Putting People First for Organizational Success}.
Academy of Management Executive, 13(2), 37-48.

\bibitem{shrm2022}
Society for Human Resource Management. (2022).
\textit{2022 Retention Report: Talent Fully Engaged?}
SHRM Research, Washington, DC.

\bibitem{gallup2022}
Gallup, Inc. (2022).
\textit{State of the Global Workplace: 2022 Report}.
Gallup Press, Washington, DC.

\bibitem{mckinsey2022}
McKinsey \& Company. (2022).
\textit{The future of work after COVID-19}.
McKinsey Global Institute Report.

\bibitem{chambers2015}
Chambers, E. G., et al. (2015).
\textit{The War for Talent}.
Harvard Business Review, 76(5), 103-108.

\bibitem{scikit2011}
Pedregosa, F., et al. (2011).
\textit{Scikit-learn: Machine Learning in Python}.
Journal of Machine Learning Research, 12, 2825-2830.

\bibitem{breiman2001}
Breiman, L. (2001).
\textit{Random Forests}.
Machine Learning, 45(1), 5-32.

\end{thebibliography}

% ============================================================================
% APÉNDICES
% ============================================================================
\newpage
\section*{Apéndice A: Fórmulas y Métricas}

\subsection*{Tasa de Rotación}
\[
\text{Rotación} = \frac{\text{Empleados que se van}}{\text{Promedio empleados en período}} \times 100
\]

\subsection*{Costo de Rotación Individual}
\[
\text{Costo} = \text{Salario} + (0.5 \times \text{Salario}) + \text{Capacitación}
\]

Para este análisis se estimó en USD \$3,000 por persona (promedio industria).

\subsection*{ROC-AUC (Área Bajo la Curva)}
\[
\text{AUC} = \int_{0}^{1} \text{TPR}(x) \, d\text{FPR}(x)
\]

Donde TPR = Tasa de Verdaderos Positivos, FPR = Tasa de Falsos Positivos

\newpage
\section*{Apéndice B: Instrucciones de Uso del Modelo}

El modelo Random Forest entrenado puede implementarse para:

\begin{enumerate}
    \item \textbf{Predicción Mensual:} Clasificar colaboradores por riesgo
    \item \textbf{Alertas Automáticas:} Flag cuando score supera umbral
    \item \textbf{Priorización de Intervenciones:} Ranking por urgencia
    \item \textbf{Evaluación de Impacto:} Seguimiento post-intervención
\end{enumerate}

\textbf{Limitaciones:}
\begin{itemize}
    \item El modelo fue entrenado con datos de mayo 2026
    \item Cambios significativos en política laboral pueden requerir reentrenamiento
    \item La correlación no implica causalidad en los predictores
\end{itemize}

\end{document}
